\scanpage{15}

% A detached point lies below-left of the opening J; its textual status is uncertain.
$\underset{\hspace{-0.6em}\unsure{\cdot|\phantom{\cdot}}}{\text{J}}$e te signale que j'ai \unsure{réfléchi} aux Gr-champs sur $X$. Si $G$ est un Groupe sur $X$, $N$ un $G$-Module
% Isolated point in the left margin at whole-source (29,246).
\reversemarginpar\margin{$\unsure{\cdot|\phantom{\cdot}}$}\normalmarginpar
les Gr-champs sur $X$ ``épinglés par $G,N$'' forment a priori une $2$-\unsure{catégori|catégorie} et même une $2$-catégorie de Picard stricte, grâce à l'opération évident à la Baer. On trouve que le complexe \inserted{(de cochaines)} tronqué à l'échelon $2$ qui leur correspond est le tronqué
\[
\tau_{\leq 2}\bigl(\mathbb R\Gamma(B_{G}\mathbin{\mathrm{mod}}X,N)[1]\bigr)
\]
(\uline{NB} la cohomologie de $\mathbb R\Gamma(B_{G}\mathbin{\mathrm{mod}}X,N)$ commence en degré $1$) Plus \unsure{géométriquement}, un Gr-champ sur $X$ \unsure{épi-g|épi-glé|épinglé} \inserted{par $(G,N)$} est essentiellement ``la même chose'' qu'une $2$-gerbe sur $B_{G}$, liée par $N$, et \deleted{\unsure{épi|pri}} munie d'une \unsure{trivialisation} au dessus de $X\approx B_{c}=(B_{G})_{/P}$ (où $P$ est le $\overset{\unsure{\deleted{s}|\ast}}{\text{o}}$bjet de $B_{G}$ ``torseur universel sous $G$''). Ces $2$-gerbes forment en fait une $3$-catégorie
% Detached point below de in the final manuscript line; no grammatical role assigned.
$\underset{\unsure{\cdot|\phantom{\cdot}}}{\text{de}}$ Picard a priori, mais il se

% archive sheet number: 14
